\section{Methods}
\label{sec:method}

\subsection{Problem Formulation and Backbone-Conditioned Approach}
Ribonucleic acid (RNA) is a nucleic acid present in all living cells. An RNA molecule, often single-stranded, has a backbone made of alternating phosphate groups and sugar ribose. RNA has four nucleotides (also known as bases): adenine (A), guanine (G), cytosine (C), and uracil (U).

We formulate RNA inverse folding as a multi-objective optimization problem over the discrete RNA sequence space $\mathcal{S}=\{A,U,C,G\}^L$ ($L$ is the RNA sequence length), where we seek sequences that jointly maximize structural fidelity and thermodynamic stability:
\begin{equation}
s^*=\arg\max_{s\in\mathcal{S}}\left[f_{\text{struct}}(s,\mathcal{G})+f_{\text{thermo}}(s)\right],
\end{equation}
where $\mathcal{G}$ denotes the target backbone geometry, $f_{\text{struct}}$ captures structural quality metrics (e.g., pLDDT, RMSD), and $f_{\text{thermo}}$ quantifies thermodynamic favorability via free energy.

To solve this, we use a backbone-conditioned policy that generates sequences given $\mathcal{G}$. Specifically, our policy $\pi_\theta(s\mid\mathcal{G})$ follows the GNN-GVP architecture from gRNAde \citep{joshi2025grnade}, encoding the 3D backbone and producing sequences autoregressively:
\begin{equation}
\pi_\theta(s\mid\mathcal{G})=\prod_{i=1}^L \pi_\theta(s_i \mid s_{<i}, \mathcal{G}).
\end{equation}
We maintain a frozen reference policy $\pi_{\mathrm{ref}}$ initialized from the pre-trained gRNAde model to provide a stable baseline and regularize distributional drift.

Optimization proceeds via multi-round Direct Preference Optimization, which progressively shifts $\pi_\theta$ toward \emph{more optimized solutions}. After each round, the improved policy supplies higher-quality preference pairs for the next round, yielding an iterative curriculum that stabilizes training while aligning the generation distribution with the multi-objective targets above.


\subsection{Multi-Round Preference Construction}

\paragraph{Progressive Refinement Strategy.}
Our iterative approach leverages improved policies from previous rounds to generate candidate sequences in previously unexplored regions of sequence space \citep{belanger2019biological}. This addresses the exploration-exploitation trade-off inherent in RNA sequence optimization, where complex structure-function relationships create numerous local optima \citep{hofacker2010barmap,jimenez2013comprehensive}. Single-round optimization frequently converges to sequences with adequate structural similarity but suboptimal thermodynamic properties or limited designability under physiological conditions. By progressively refining discrimination criteria across rounds, our method identifies increasingly subtle quality differences that drive convergence toward globally optimal solutions, balancing structural accuracy with biophysical stability \citep{gonzalez2017preferential}. 
We operationalize this progressive refinement through round-wise tightening of preference criteria, as detailed below.

\paragraph{Preference Criteria.}
We establish high-quality preference pairs using systematic statistical thresholds tailored to the RNA inverse folding task. Given two sequences $s^w$ (winner/chosen) and $s^l$ (loser/rejected), a preference is recorded only when both filter and margin conditions are satisfied:

\begin{align}
\text{\textbf{Filter:}} \quad & \text{pLDDT}(g(s^w)) > 0.70 \land \text{RMSD}(\mathcal{G}, g(s^w)) < 8\,\text{\AA} \\
% --- Changes below ---
\text{\textbf{Margin:}} \quad & \underbrace{\left(\bigwedge_{m\in\{\text{pLDDT, RMSD}\}}|m(s^w)-m(s^l)|>\gamma_r\cdot\sigma_m\right)}_{\text{Structural Margin}} \land \underbrace{\left(\text{MFE}(s^w) < \text{MFE}(s^l)\right)}_{\text{Thermodynamic Preference}}
\end{align}


where $g$ denotes RNA folding approaches (RhoFold+ \citep{shen2024accurate} in this work). $\sigma_m$ is the standard deviation of metric $m(\cdot)$ estimated from the current candidate pool. These criteria exclude candidates unlikely to fold correctly, ensuring that preference labels are assigned within a biologically realistic design space and reducing downstream noise. The statistical margin $\gamma_r\sigma_m$ enforces the \emph{significance} of the structural preferences, effectively improving the signal-to-noise ratio of the implicit reward; normalization by $\sigma_m$ renders metrics on different scales commensurate and stabilizes the DPO likelihood-ratio updates. 



Concretely, we implement this progressive refinement of structure and energy by decreasing $\gamma_r$ across iterations. Early rounds use looser margins to rapidly establish broad foldability; later rounds tighten the margins to capture subtle, near–native improvements (e.g., local geometry). This round-wise annealing functions as Curriculum Learning (CL): large early gaps yield high-confidence preferences that reduce label noise and reward sparsity, while smaller later gaps provide denser, finer-grained gradient signals. The schedule mirrors an exploration-to-exploitation transition and implicitly enforces a trust-region constraint relative to the reference policy, stabilizing updates of large autoregressive models.  


\subsection{Reinforcement Learning from Physical Feedback via DPO}
\label{subsec:rlpf_dpo}
We construct preference pairs $(s^w, s^l)$ by sampling candidate sequences from the current policy $\pi_\theta(\,\cdot\,\mid\mathcal{G})$ for each target RNA backbone $\mathcal{G}$ and then applying the filter and margin criteria described above. This procedure ensures that $s^w$ and $s^l$ reflect biologically realistic, thermodynamically feasible regions of sequence space.



Our objective maximizes the preference likelihood ratio between $s^w$ and $s^l$ while anchoring against the reference policy:
\begin{equation}
\mathcal{L}_{\mathrm{DPO}}
= -\mathbb{E}\!\left[
\log \sigma \!\Bigl(
\beta \log \frac{\pi_\theta(s^w \mid \mathcal{G})}{\pi_\theta(s^l \mid \mathcal{G})}
- \beta \log \frac{\pi_{\mathrm{ref}}(s^w \mid \mathcal{G})}{\pi_{\mathrm{ref}}(s^l \mid \mathcal{G})}
\Bigr)
\right].
\end{equation}
where it directly shapes the policy toward generating sequences with better structural fidelity and thermodynamic stability. 
while the reference-policy term acts as a trust region to prevent destabilizing updates, preventing regressions toward suboptimal folds.

To stabilize training further, we add supervised fine-tuning (SFT) loss on the preferred sequences:
\begin{equation}
\mathcal{L}
= \mathcal{L}_{\mathrm{DPO}}
+ \lambda_{\mathrm{SFT}}\,
\mathbb{E}_{s^w}\bigl[-\log \pi_\theta(s^w \mid \mathcal{G})\bigr],
\end{equation}
where $\lambda_{\mathrm{SFT}}$ is a hyperparameter controlling the strength of this term.  
The SFT component acts as an explicit maximum-likelihood ``anchor" on sequences with higher designability and thermostability potential, reducing catastrophic drift and mode collapse often seen in RL-based fine-tuning.  

% \textcolor{blue}{Mention designability in the introduction section}

% \subsection{SSTT Benchmark: Comprehensive Evaluation Framework}
% \label{subsec:sstt}

% Unlike existing RNA inverse folding benchmarks, which mainly evaluate sequence recovery and RMSD, we introduce the SSTT benchmark to assess performance across four complementary dimensions. This framework integrates thermodynamic stability and ensemble properties alongside traditional structural metrics, offering a multidimensional view of design quality. Detailed calculation procedures are provided in Appendix~\ref{subsec:sstt_detail}. {\textbf{S}equence Axis} captures how closely designed sequences resemble native ones and how broadly they explore sequence space. We assess recovery rate \citep{dauparas2022proteinmpnn} to measure native residue similarity, and 3-mer diversity \citep{shannon1948mathematical,Bokulich2024diversitykmer} to quantify sequence exploration capacity. {\textbf{S}econdary Structure Axis} evaluates whether designed sequences fold back into their intended secondary structures. We measure forward-folding self-consistency using Matthews Correlation Coefficient \citep{matthews1975comparison,mccmoreaccurate} on EternaFold \citep{waymentsteele2022eternafold} predictions, capturing base-pairing fidelity. {\textbf{T}ertiary Structure Axis} assesses three-dimensional fidelity at both global and local levels as well as steric and torsional realism. We evaluate TM-score \citep{zhang2004tmscore}, GDT, and RMSD for global similarity; pLDDT \citep{jumper2021alphafold2} for confidence; lDDT for local structure preservation; interaction network fidelity \citep{parisien2009newmetricsinf} for base-pairing accuracy; clash score \citep{word1999clashscore,davis2007molprobity} for geometric feasibility; and backbone torsion quality \citep{ramachandran1963stereochemistry} for stereochemical realism. {\textbf{T}hermostability Axis} quantifies the thermodynamic robustness of designed sequences under physiological conditions. We incorporate minimum free energy \citep{zuker1981optimal} for stability, ensemble defect \citep{zadeh2011nucleic} to measure target structure specificity, target structure probability \citep{McCaskill1990TheEquilibrium} to quantify conformational preference, Shannon entropy for ensemble sharpness, and melting temperature for thermal resilience. 
 
% % comprehensively


\subsection{SSTT Benchmark: Comprehensive Evaluation Framework}
\label{subsec:sstt}

Unlike existing RNA inverse folding benchmarks, which mainly evaluate sequence recovery and RMSD, we introduce the SSTT benchmark to assess performance across four complementary dimensions. This framework integrates thermodynamic stability and ensemble properties alongside traditional structural metrics, offering a multidimensional view of design quality. Detailed calculation procedures are provided in Appendix~\ref{subsec:sstt_detail}. 
{\textbf{S}equence Axis} captures how closely designed sequences resemble native ones and how broadly they explore sequence space. We assess recovery rate \citep{dauparas2022proteinmpnn} to measure native residue similarity, and 3-mer diversity \citep{shannon1948mathematical,Bokulich2024diversitykmer} to quantify sequence exploration capacity.
{\textbf{S}econdary Structure Axis} evaluates whether designed sequences fold back into their intended secondary structures. We measure forward-folding self-consistency using Matthews Correlation Coefficient \citep{matthews1975comparison,mccmoreaccurate} on EternaFold \citep{waymentsteele2022eternafold} predictions, capturing base-pairing fidelity. 
{\textbf{T}ertiary Structure Axis} assesses three-dimensional fidelity at both global and local levels as well as steric realism. We evaluate TM-score \citep{zhang2004tmscore}, GDT, and RMSD for global similarity; pLDDT \citep{jumper2021alphafold2} for confidence; { interaction network fidelity \citep{parisien2009newmetricsinf} for base-pairing accuracy; and clash score \citep{word1999clashscore,davis2007molprobity} for geometric feasibility.} 
{\textbf{T}hermostability Axis} quantifies the thermodynamic robustness of designed sequences under physiological conditions. We incorporate minimum free energy \citep{zuker1981optimal} for stability, { and ensemble defect \citep{zadeh2011nucleic} to measure target structure specificity.}